\documentclass[letterpaper]{article}
\usepackage[submission]{aaai2026}

\usepackage{amsmath,amssymb,amsfonts}
\usepackage{booktabs}
\usepackage{graphicx}
\usepackage{xcolor}
\usepackage{algorithm}
\usepackage{algorithmic}
\usepackage{multirow}
\usepackage{url}

\title{FinVerse: A Hierarchical Market-State World Model as the Cognitive Core of Financial Agents}

\author{
Anonymous Authors
}

\affil{
Paper under review
}

\begin{document}
\maketitle

\begin{abstract}
Financial decision making requires an agent to reason about a partially observed, noisy, and regime-shifting world. Rather than treating market prediction as a direct supervised forecasting problem, we study the construction of a financial world model that can serve as the cognitive core of a future autonomous financial agent. We propose FinVerse, a hierarchical market-state world model centered on a Hierarchical Market State Constructor (HMSC). FinVerse tokenizes temporal price dynamics and cross-sectional market context through dual vector quantization, aligns continuous news, macro, and neighbor-aware market features into a shared representation space, and learns recurrent probabilistic latent dynamics for future market imagination. This design supports counterfactual questions of the form: if an agent follows a candidate strategy under the current market state, how may future returns and risks change? Experiments on U.S. market data from 2020--2025 compare FinVerse with LSTM, GRU, Transformer, PatchTST, Kronos-mini, and Vanilla RSSM baselines. Results show that FinVerse obtains the best short-horizon forecasting error under the same evaluation protocol, while ablation studies indicate that discrete market-state tokenization, cross-asset context, and probabilistic dynamics contribute complementary benefits. We present FinVerse as a step toward sample-efficient financial agents that improve decision quality through internal market simulation before real-world interaction.
\end{abstract}

\section{Introduction}
Modern financial markets are difficult environments for autonomous agents. The agent observes incomplete and heterogeneous signals, including prices, volume, macro indicators, firm events, and cross-asset dependencies. These signals are non-stationary and often change meaning under different market regimes. A useful financial agent therefore requires more than a point predictor. It needs an internal model of the financial world: a compact representation of the current market state, a mechanism for imagining plausible future states, and an evaluation mechanism for comparing candidate strategies before execution.

This paper focuses on this internal model. We do not attempt to build a complete end-to-end trading agent. Instead, we study the cognitive core that such an agent would rely on. The central question is: how can an agent construct a market state representation that supports forecasting, imagination, and strategy evaluation under limited real-world interaction?

We propose FinVerse, a Hierarchical Market-State World Model. The core module is the Hierarchical Market State Constructor (HMSC), which maps multi-modal market observations into a unified latent state. HMSC uses dual vector-quantized tokenization to discretize temporal price patterns and cross-sectional market context, while continuous news, macro, and neighbor-aware market features are projected into the same representation space. A recurrent probabilistic world model then rolls the latent state forward and predicts future return trajectories.

Our contributions are threefold:
\begin{itemize}
    \item We formulate financial world modeling as the cognitive core of an agent, emphasizing market-state construction, imagination, and strategy evaluation rather than direct one-step prediction alone.
    \item We introduce HMSC, a hierarchical multi-modal state constructor with dual VQ tokenization for temporal and cross-sectional market patterns.
    \item We evaluate FinVerse against mainstream forecasting models and report rollout fidelity diagnostics showing how imagined return trajectories match realized future trajectories.
\end{itemize}

\section{Related Work}
\paragraph{Financial Forecasting.}
Financial forecasting has been studied with recurrent networks, temporal convolutional models, and Transformer-based architectures. LSTM and GRU models are widely used for sequential return prediction, while PatchTST and related time-series Transformers improve long-horizon modeling by segmenting time series into patches. These methods are strong predictive baselines, but they usually learn direct mappings from historical observations to future targets.

\paragraph{World Models and Imagination.}
World models learn latent dynamics that allow an agent to simulate future trajectories before acting. In reinforcement learning, models such as RSSM and Dreamer learn compact latent states and optimize policies through imagined rollouts. FinVerse adapts this philosophy to financial markets, where the imagined future must represent stochastic returns, cross-asset interactions, and regime-sensitive market states.

\paragraph{Financial Agents.}
Existing financial-agent systems often focus on portfolio allocation, execution, or language-driven decision making. Our goal is narrower and more foundational: we design the internal market-state world model that a later full financial agent can use for reasoning.

\section{Methodology}
\subsection{Problem Formulation}
At each time step $t$, the agent observes a multi-modal market observation
\begin{equation}
    o_t = \{x_t^{p}, x_t^{n}, x_t^{m}, G_t\},
\end{equation}
where $x_t^{p}$ denotes historical price and volume features, $x_t^{n}$ denotes news or event-derived features, $x_t^{m}$ denotes macro features, and $G_t$ denotes the cross-asset neighborhood context. Given a candidate action or strategy descriptor $a_t$, the objective is to learn a world model that estimates future return-related outcomes:
\begin{equation}
    p_\theta(o_{t+1:t+H}, r_{t+1:t+H} \mid o_{\leq t}, a_{t:t+H-1}).
\end{equation}

\subsection{Hierarchical Market State Constructor}
The HMSC maps heterogeneous observations into a shared latent market state:
\begin{equation}
    s_t = \phi_{\mathrm{HMSC}}(x_t^{p}, x_t^{n}, x_t^{m}, G_t).
\end{equation}
The design separates discrete market-state tokenization from continuous feature projection.

\paragraph{Temporal Price Tokenization.}
Given a lookback window $X_t^p \in \mathbb{R}^{L \times d_p}$, an encoder produces a temporal representation:
\begin{equation}
    h_t^{\mathrm{temp}} = f_{\mathrm{temp}}(X_t^p).
\end{equation}
A vector quantizer maps it to the nearest codebook vector:
\begin{equation}
    k_t^{\mathrm{temp}} = \arg\min_k \|h_t^{\mathrm{temp}} - e_k^{\mathrm{temp}}\|_2^2,
    \quad
    z_t^{\mathrm{temp}} = e_{k_t^{\mathrm{temp}}}^{\mathrm{temp}}.
\end{equation}

\paragraph{Cross-sectional Tokenization.}
For the same trading day, cross-asset structure is encoded as
\begin{equation}
    h_t^{\mathrm{cross}} = f_{\mathrm{cross}}(X_t^p, G_t),
\end{equation}
and quantized by
\begin{equation}
    k_t^{\mathrm{cross}} = \arg\min_k \|h_t^{\mathrm{cross}} - e_k^{\mathrm{cross}}\|_2^2,
    \quad
    z_t^{\mathrm{cross}} = e_{k_t^{\mathrm{cross}}}^{\mathrm{cross}}.
\end{equation}
The VQ objective is
\begin{equation}
    \mathcal{L}_{\mathrm{VQ}} =
    \| \mathrm{sg}[h_t] - z_t \|_2^2
    + \beta \| h_t - \mathrm{sg}[z_t] \|_2^2,
\end{equation}
where $\mathrm{sg}[\cdot]$ denotes stop-gradient.

\paragraph{Continuous Tokens and Shared Space.}
News, macro, and cross-asset context features are projected into the same latent space:
\begin{equation}
    \tilde{h}_t^n = \psi_n(x_t^n), \quad
    \tilde{h}_t^m = \psi_m(x_t^m), \quad
    \tilde{h}_t^g = \psi_g(G_t).
\end{equation}
In the current implementation, $G_t$ is represented by a target asset and its nearest cross-asset neighbors. The graph-context branch encodes the neighbor price histories and pools them into $\tilde{h}_t^g$. We therefore use the term graph-context encoding rather than claiming a full message-passing GNN.
The unified market state is obtained by fusion:
\begin{equation}
    s_t = \phi_{\mathrm{fuse}}\left(
    [z_t^{\mathrm{temp}}, z_t^{\mathrm{cross}}, \tilde{h}_t^n, \tilde{h}_t^m, \tilde{h}_t^g]
    \right).
\end{equation}

\subsection{Recurrent Financial World Model}
FinVerse learns a recurrent latent dynamics model. The posterior state is
\begin{equation}
    q_\phi(z_t \mid s_t) = \mathcal{N}(\mu_\phi(s_t), \sigma_\phi(s_t)),
\end{equation}
and the transition prior is
\begin{equation}
    p_\theta(z_{t+1} \mid z_t, a_t) =
    \mathcal{N}(\mu_\theta(h_{t+1}), \sigma_\theta(h_{t+1})),
\end{equation}
where
\begin{equation}
    h_{t+1} = \mathrm{GRUCell}([z_t, a_t], h_t).
\end{equation}
The decoder predicts future returns:
\begin{equation}
    \hat{y}_{t+h} = g_\theta(z_{t+h}), \quad h=1,\ldots,H.
\end{equation}
We also attach an auxiliary regime head
\begin{equation}
    \hat{c}_t = \arg\max \; g_{\mathrm{reg}}(z_t),
\end{equation}
where $c_t \in \{\mathrm{bear}, \mathrm{sideway}, \mathrm{bull}\}$ is derived from the average realized return over the first five future steps. In our experimental protocol, bear and bull labels are assigned when this short-horizon realized return is below $-1\%$ or above $+1\%$, respectively; all remaining cases are sideway.

\subsection{Strategy Imagination}
For a candidate strategy $\pi$, the model imagines a latent rollout:
\begin{equation}
    z_{t+h}^{\pi} \sim p_\theta(z_{t+h} \mid z_{t+h-1}^{\pi}, a_{t+h-1}^{\pi}).
\end{equation}
The expected value of the strategy can be estimated by
\begin{equation}
    J(\pi) = \mathbb{E}\left[\sum_{h=1}^{H}\gamma^{h-1}\hat{r}_{t+h}^{\pi}\right]
    - \lambda \widehat{\mathrm{Risk}}(\pi).
\end{equation}
This formulation supports the cognitive question: if the agent follows this strategy, how may its future return and risk change?

\subsection{Counterfactual Diagnostic Protocol}
To make the world-model claim testable, we define a controlled counterfactual diagnostic. Given a factual observation $o_t$ and a perturbed macro context $\tilde{o}_t^{m}=o_t^{m}+\delta$, the model produces two imagined trajectories:
\begin{equation}
    \hat{y}_{t+1:t+H}=F_\theta(o_t,a_{t:t+H-1}),
    \quad
    \tilde{y}_{t+1:t+H}=F_\theta(\tilde{o}_t,a_{t:t+H-1}).
\end{equation}
We measure shock sensitivity by the mean absolute prediction change and direction-flip rate:
\begin{equation}
    \Delta_{\mathrm{cf}} =
    \frac{1}{H}\sum_{h=1}^{H}
    |\tilde{y}_{t+h}-\hat{y}_{t+h}|.
\end{equation}
This diagnostic is used to test whether the learned latent world reacts to controlled market-context changes. We keep it as a diagnostic protocol unless retrained checkpoints show stable responses.

\subsection{Crisis-Window Simulation Diagnostic}
We retain crisis simulation as a stress-window diagnostic rather than a broad claim of fully causal crisis generation. A crisis window is defined as a test sample whose average realized return over the first five future steps falls below a fixed stress threshold $\tau_c$:
\begin{equation}
    \frac{1}{5}\sum_{h=1}^{5} y_{t+h} \leq \tau_c,
    \quad \tau_c=-0.02.
\end{equation}
On this subset, we evaluate crisis-window rollout fidelity:
\begin{equation}
    \mathrm{CrisisStateMSE}@h =
    \frac{1}{|\mathcal{C}|}\sum_{i\in\mathcal{C}}
    \|\hat{y}_{i,t+h}-y_{i,t+h}\|_2^2.
\end{equation}
This test asks whether imagined trajectories remain accurate under realized market stress. It does not claim to recover all causal mechanisms behind historical events such as COVID or rate-hike episodes.

\subsection{Training Objective}
The total training loss combines prediction, dynamics regularization, and tokenization:
\begin{equation}
    \mathcal{L} =
    \mathcal{L}_{\mathrm{pred}}
    + \alpha \mathcal{L}_{\mathrm{KL}}
    + \eta \mathcal{L}_{\mathrm{VQ}}
    + \rho \mathcal{L}_{\mathrm{reg}},
\end{equation}
where
\begin{equation}
    \mathcal{L}_{\mathrm{pred}} =
    \frac{1}{H}\sum_{h=1}^{H}
    \|\hat{y}_{t+h} - y_{t+h}\|_2^2,
\end{equation}
and
\begin{equation}
    \mathcal{L}_{\mathrm{KL}} =
    D_{\mathrm{KL}}(q_\phi(z_t \mid s_t) \| p_\theta(z_t \mid z_{t-1}, a_{t-1})).
\end{equation}
The auxiliary regime loss $\mathcal{L}_{\mathrm{reg}}$ is a cross-entropy loss over bear, sideway, and bull labels derived from short-horizon realized returns. It is used as an auxiliary diagnostic signal rather than the primary training target.

\begin{algorithm}[t]
\caption{FinVerse World Model Learning and Strategy Evaluation}
\label{alg:finverse}
\begin{algorithmic}[1]
\STATE Observe multi-modal market input $o_t=\{x_t^p,x_t^n,x_t^m,G_t\}$.
\STATE Construct temporal and cross-sectional VQ tokens.
\STATE Project news, macro, and cross-asset context features into the shared latent space.
\STATE Fuse all tokens into hierarchical market state $s_t$.
\STATE Infer posterior latent state $z_t \sim q_\phi(z_t \mid s_t)$.
\STATE Roll forward recurrent dynamics under candidate actions.
\STATE Decode imagined returns, risks, and auxiliary market regime.
\STATE Select or evaluate strategies according to risk-adjusted expected value.
\end{algorithmic}
\end{algorithm}

\section{Experiments}
\subsection{Data}
We evaluate on U.S. market data from 2020--2025. The processed dataset contains 90 actively used assets and ETFs, daily price sequences, price-derived macro proxies, event/news proxy features, and cross-asset neighborhood information. Each sample uses a 30-day lookback window and predicts future returns up to 30 days. We use chronological train, validation, and test splits. The present version does not use chart images or an explicit memory-retrieval bank; these are left for future extensions to avoid unsupported modality claims.

\subsection{Experimental Protocol}
We fix a single protocol for all model variants to avoid post-hoc comparison. Unless otherwise specified, models use hidden dimension 128, latent dimension 128, return targets, batch size 128, seed 42, KL weight 0.001, VQ weight 0.001, and regime weight 0.1 for regime-supervised runs. The short protocol uses 1,000 training episodes and 300 validation episodes; the medium protocol uses 3,000 training episodes and 800 validation episodes; the large protocol uses 6,000 training episodes and 1,200 validation episodes. All reported comparisons should use the same chronological test split.

We define five experiment groups. First, forecasting comparison evaluates FinVerse against mainstream sequence models. Second, rollout fidelity compares imagined trajectories with realized future trajectories to test world-model behavior. Third, regime prediction evaluates the auxiliary bear/sideway/bull head using accuracy and macro-F1 after retraining with regime supervision. Fourth, crisis-window simulation evaluates rollout fidelity on realized stress windows. Fifth, ablation studies remove Dual VQ, cross-asset context, probabilistic dynamics, and non-price modalities under the same protocol.

\subsection{Baselines}
The full baseline pool includes LSTM, GRU, Transformer, PatchTST, TimesFM-style, Chronos-mini, Kronos-mini, Vanilla RSSM, and Dreamer-style RSSM. TimesFM-style and Chronos-mini are lightweight train-from-scratch baselines that follow the patch-token and discretized-token design principles of the corresponding foundation models; they are not official pretrained checkpoints. Dreamer-style RSSM is a compact recurrent world-model baseline with stochastic latent dynamics. The newly added TimesFM-style, Chronos-mini, and Dreamer-style RSSM rows are configured for the next protocol run and should be populated before final submission. BUY\&HOLD is retained as a sanity-check investment baseline in auxiliary diagnostics, but portfolio PnL is not used as the main evidence because short-window simulated returns are highly sensitive to annualization and ranking noise.

\subsection{Metrics}
Forecasting quality is measured by MSE and MAE at horizons $1,5,10,20,30$. Cross-sectional predictive quality is measured by IC and RankIC:
\begin{equation}
    \mathrm{IC}_h = \mathrm{corr}(\hat{y}_{t+h}, y_{t+h}),
\end{equation}
\begin{equation}
    \mathrm{RankIC}_h = \mathrm{corr}(\mathrm{rank}(\hat{y}_{t+h}), \mathrm{rank}(y_{t+h})).
\end{equation}
Volatility quality is measured by the MAE between predicted and realized path volatility. We intentionally exclude annualized portfolio PnL from the main comparison table because the current short-window top-k portfolios can produce unstable daily returns that are not suitable as primary evidence for a world-model paper.
To directly evaluate world-model behavior, we additionally report rollout fidelity. Given an imagined trajectory $\hat{y}_{t+1:t+H}$ and the realized future trajectory $y_{t+1:t+H}$, we compute
\begin{equation}
    \mathrm{StateMSE}@h =
    \frac{1}{N}\sum_{i=1}^{N}
    \|\hat{y}_{i,t+h} - y_{i,t+h}\|_2^2.
\end{equation}
This metric measures whether the latent rollout remains close to the realized market trajectory over multiple horizons.
For regime diagnostics, we report regime accuracy and macro-F1. These metrics are only reported for checkpoints trained with the regime-supervised objective. For counterfactual diagnostics, we report the mean absolute prediction change and the direction-flip rate under a controlled macro-context perturbation. For crisis-window simulation, we report CrisisStateMSE on stress-window samples whose realized short-horizon returns fall below the crisis threshold.

\begin{table*}[t]
\centering
\caption{Main forecasting comparison on 500 held-out test episodes. Lower values indicate better prediction accuracy.}
\label{tab:main_forecasting}
\resizebox{\textwidth}{!}{
\begin{tabular}{lcccccccccc}
\toprule
Model & MSE@1 & MSE@5 & MSE@10 & MSE@20 & MSE@30 & MAE@1 & MAE@5 & MAE@10 & MAE@20 & MAE@30 \\
\midrule
LSTM & 0.0011 & 0.0078 & \textbf{0.0064} & 0.0171 & 0.0232 & 0.0256 & 0.0753 & \textbf{0.0626} & 0.1050 & 0.1284 \\
PatchTST & 0.0012 & 0.0068 & 0.0336 & 0.0136 & 0.0625 & 0.0278 & 0.0671 & 0.1677 & 0.1000 & 0.2025 \\
Kronos-mini & 0.0028 & 0.0118 & 0.0092 & 0.0209 & 0.0332 & 0.0484 & 0.0967 & 0.0760 & 0.1191 & 0.1554 \\
Vanilla RSSM & 0.0503 & \textbf{0.0033} & 0.0064 & 0.0133 & \textbf{0.0231} & 0.2230 & \textbf{0.0441} & 0.0629 & 0.0987 & 0.1286 \\
\textbf{FinVerse} & \textbf{0.0008} & 0.0044 & 0.0064 & 0.0135 & 0.0235 & \textbf{0.0202} & 0.0507 & 0.0628 & 0.0979 & \textbf{0.1282} \\
GRU & 0.0011 & 0.0034 & 0.0077 & 0.0134 & 0.0260 & 0.0252 & 0.0447 & 0.0692 & 0.0994 & 0.1326 \\
Transformer & 0.0019 & 0.0154 & 0.0232 & \textbf{0.0129} & 0.0232 & 0.0386 & 0.1113 & 0.1323 & \textbf{0.0968} & 0.1303 \\
\bottomrule
\end{tabular}
}
\end{table*}

\begin{table*}[t]
\centering
\caption{Rollout fidelity for imagined return trajectories. Lower StateMSE indicates that the imagined trajectory better matches the realized future trajectory.}
\label{tab:rollout_fidelity}
\resizebox{\textwidth}{!}{
\begin{tabular}{lccccc}
\toprule
Model & StateMSE@1 & StateMSE@5 & StateMSE@10 & StateMSE@20 & StateMSE@30 \\
\midrule
Full FinVerse & 0.0003 & 0.0015 & 0.0027 & 0.0046 & \textbf{0.0071} \\
w/o Dual VQ & 0.0004 & 0.0015 & 0.0027 & 0.0046 & 0.0072 \\
w/o Cross-Asset Ctx & 0.0005 & \textbf{0.0015} & \textbf{0.0027} & 0.0046 & 0.0071 \\
w/o Probabilistic WM & \textbf{0.0003} & 0.0016 & 0.0027 & 0.0046 & 0.0072 \\
Price Only & 0.0008 & 0.0015 & 0.0030 & \textbf{0.0045} & 0.0071 \\
\bottomrule
\end{tabular}
}
\end{table*}

\begin{table*}[t]
\centering
\caption{Financial prediction diagnostics on 100 complete test trading days. We report forecasting error, cross-sectional IC/RankIC, and volatility error, and exclude portfolio PnL from the main table because short-window daily portfolio returns are highly sensitive to ranking noise and annualization choices.}
\label{tab:financial_metrics}
\resizebox{\textwidth}{!}{
\begin{tabular}{lccccc}
\toprule
Model & MSE@1 & MSE@30 & IC & RankIC & Vol. MAE \\
\midrule
LSTM & \textbf{0.0442} & 0.0662 & -0.0083 & -0.0219 & 0.0284 \\
PatchTST & 0.0802 & 0.1395 & 0.0481 & 0.0195 & 0.1303 \\
Kronos-mini & 0.0972 & 0.1545 & 0.0580 & 0.0406 & 0.1316 \\
Vanilla RSSM & 0.0694 & 0.0662 & -0.0046 & 0.0104 & 0.0550 \\
FinVerse-Small & 0.0451 & 0.0646 & \textbf{0.0587} & \textbf{0.0657} & 0.0336 \\
FinVerse-Medium & 0.0459 & 0.0660 & 0.0226 & 0.0249 & 0.0387 \\
FinVerse-Large & 0.0458 & 0.0656 & -0.0846 & -0.0502 & 0.0397 \\
GRU & 0.0457 & \textbf{0.0622} & -0.1674 & -0.0800 & \textbf{0.0276} \\
Transformer & 0.0544 & 0.0971 & 0.0298 & 0.0512 & 0.0958 \\
\bottomrule
\end{tabular}
}
\end{table*}

\begin{table*}[t]
\centering
\caption{Same-protocol ablation results for forecasting quality.}
\label{tab:ablation_forecast}
\resizebox{\textwidth}{!}{
\begin{tabular}{lcccccccc}
\toprule
Model & MSE@1 & MSE@5 & MSE@30 & MAE@1 & MAE@30 & IC & RankIC & Vol. MAE \\
\midrule
Full FinVerse & \textbf{0.0008} & 0.0044 & 0.0235 & \textbf{0.0202} & \textbf{0.1282} & -0.0659 & -0.0670 & 0.0471 \\
w/o Dual VQ & 0.0503 & \textbf{0.0033} & \textbf{0.0231} & 0.2230 & 0.1286 & -0.1139 & -0.1366 & 0.0388 \\
w/o Cross-Asset Ctx & 0.0011 & 0.0041 & 0.0231 & 0.0259 & 0.1284 & -0.0626 & -0.0593 & 0.0404 \\
w/o Probabilistic WM & 0.0039 & 0.0045 & 0.0236 & 0.0585 & 0.1282 & 0.0854 & 0.0840 & \textbf{0.0296} \\
Price Only & 0.0064 & 0.0056 & 0.0240 & 0.0762 & 0.1329 & \textbf{0.2227} & \textbf{0.2422} & 0.0470 \\
\bottomrule
\end{tabular}
}
\vspace{2pt}
\footnotesize{All variants use the same training split, validation split, seed, return target, and 500 held-out test episodes. Lower is better for MSE, MAE, and Vol. MAE; higher is better for IC and RankIC.}
\end{table*}


\subsection{Main Results}
Table~\ref{tab:main_forecasting} shows that FinVerse achieves the best one-step forecasting performance among the compared models, with MSE@1 of 0.0008 and MAE@1 of 0.0202. This result supports the role of HMSC in constructing an informative short-horizon market state. Vanilla RSSM performs competitively at several longer horizons, indicating that recurrent latent dynamics are useful, but its weak one-step error suggests that dynamics alone are insufficient without the proposed market-state constructor.

Table~\ref{tab:rollout_fidelity} evaluates whether the learned world model can roll out imagined return trajectories that remain close to realized future trajectories. Full FinVerse obtains the best short-horizon rollout fidelity at StateMSE@1 and remains competitive at longer horizons. Some ablated variants obtain slightly lower long-horizon MSE, which suggests that the current rollout objective still favors smooth long-horizon averages. We therefore interpret rollout fidelity as direct world-model evidence, while treating regime-supervised and crisis-window simulation results as additional diagnostics to be filled after retraining.

Table~\ref{tab:financial_metrics} reports financial prediction diagnostics, including cross-sectional IC, RankIC, and volatility error. FinVerse-Small obtains the strongest IC and RankIC among the compared models, suggesting that HMSC improves cross-sectional signal quality even when we avoid unstable short-window portfolio PnL as a main result.

\subsection{Ablation Study}
Table~\ref{tab:ablation_forecast} evaluates the contribution of each module under the same forecasting protocol. Removing Dual VQ severely increases one-step forecasting error, suggesting that discrete temporal and cross-sectional market tokens are important for immediate state construction. Removing the cross-asset context branch keeps pointwise forecasting close but weakens cross-sectional ranking behavior, indicating that neighbor-aware market context matters beyond pointwise error alone. Removing probabilistic world-model dynamics worsens short-horizon trajectory fidelity, supporting the importance of imagined latent transitions.

\section{Limitations}
This paper studies the world-model component of a financial agent, not a complete autonomous trading system. The current experiments use a short training protocol for fast comparison, and portfolio metrics are noisy under limited test windows. News features are currently represented by public event proxies rather than full-scale professional news embeddings. The current model also does not include chart-image encoding or an explicit episodic memory-retrieval module. Crisis simulation is formulated as crisis-window rollout fidelity rather than full causal reconstruction of historical crises. Future work should extend the dataset, conduct multi-seed and multi-market evaluation, strengthen supervised regime and crisis-window diagnostics, and integrate a full actor-critic policy trained in imagined market rollouts.

\section{Conclusion}
We introduced FinVerse, a hierarchical market-state world model designed as the cognitive core of financial agents. By combining dual VQ tokenization, multi-modal state fusion, neighbor-aware cross-asset context, and recurrent probabilistic dynamics, FinVerse enables an agent to encode the current market, imagine future trajectories, and evaluate candidate strategies internally. Experiments show strong short-horizon forecasting performance and provide ablation evidence for the complementary roles of tokenization, cross-asset context, and probabilistic dynamics. FinVerse provides a foundation for future financial agents that learn more efficiently by thinking through simulated market futures before acting in the real world.

\end{document}
